\documentclass[11pt,a4paper]{letter}
\usepackage[utf8]{inputenc}
\usepackage[T1]{fontenc}
\usepackage[french]{babel}
\usepackage[left=2.5cm,right=2.5cm,top=2cm,bottom=2cm]{geometry}
\usepackage{hyperref}
\usepackage{xcolor}

% --- Couleurs Worldline ---
\definecolor{primary}{RGB}{0, 70, 135}

\hypersetup{colorlinks=true, urlcolor=primary}

\signature{Loïc Aron MBASSI EWOLO}
\address{Loïc Aron MBASSI EWOLO\\
Cergy, France\\
(+33) 7 59 52 33 98\\
\href{mailto:loicaron.mbassiewolo@ensea.fr}{loicaron.mbassiewolo@ensea.fr}\\
\href{https://www.linkedin.com/in/loic-mbassi}{LinkedIn} | \href{https://github.com/Nameless0l}{GitHub}}

\begin{document}

\begin{letter}{Worldline\\
Seclin}

\opening{Madame, Monsieur,}

Élève-ingénieur en 2A à l'ENSEA (double diplôme avec Polytechnique Yaoundé), je candidate au stage de développement fullstack Python pour le portail d'observabilité.

Le sujet me parle : mixer \textbf{Python}, \textbf{IA} et développement web, c'est exactement ce que j'aime faire. Au MILA cet été, j'ai développé des pipelines PyTorch pour analyser le comportement de modèles -- pas très loin de ce que vous voulez faire pour observer vos services en production.

Côté ML, j'ai de l'expérience avec les modèles de classification (Random Forest, SVM, XGBoost) grâce au hackathon IndabaX où j'ai fini 2e. Les modèles de séries temporelles (LSTM, Prophet, ARIMA) mentionnés dans l'offre, je les connais en théorie via mes cours, et j'ai envie de les appliquer sur des données réelles.

Pour le fullstack, j'ai 2 ans d'expérience sur une plateforme de télémédecine (SOSDOCTA) : API REST, dashboard, système de paiement. J'ai aussi travaillé sur une architecture microservices avec Docker et RabbitMQ.

J'utilise déjà des assistants IA pour coder (GitHub Copilot), donc ça ne me fait pas peur d'intégrer ça dans mon workflow.

Le côté "concevoir et challenger les solutions" m'intéresse -- j'aime comprendre les besoins avant de coder. Et le secteur des paiements, c'est un domaine où la fiabilité et l'observabilité sont critiques.

Disponible pour un entretien.

\closing{Cordialement,}

\end{letter}
\end{document}
